A retrieval system must decide how much text to return as one passage.
A question asking for a dataset name may need only a short phrase,
while an explanation may require several connected statements. A longer
passage can preserve useful context, but it can also bury the relevant
detail in unnecessary material.

This thesis asks whether adapting chunk size to a question or a region
of a paper recovers evidence more effectively than using one size
throughout. The study focuses on the retrieval stage of
retrieval-augmented generation (RAG). It evaluates evidence selection;
generated answers are outside its scope.

The experimental environment was built around QASPER, a dataset of
questions about scientific papers. Each paper was represented at five
granularities: 10, 20, 40, 80, and 160 GPT-2 tokens, using fixed
\texttt{text-embedding-3-small} embeddings stored in Qdrant. Retrieval
was restricted to the associated paper, isolating evidence selection
from document discovery.
The frozen question-level population contains 2,245 training questions
and 924 validation questions. The main retrieval protocol returns five
chunks per question. Its primary metric, joined evidence-token F1,
measures lexical overlap between the combined retrieved text and the
annotated evidence.

The work began with a Weaviate prototype before moving to Qdrant.
Resumable ingestion and checks of identifiers, offsets, and collection
completeness established a consistent experimental pipeline.

Fixed-size retrieval provided the main reference. Among the five sizes,
40 tokens achieved the highest mean joined retrieval F1, 0.3159. The
offline best-size oracle achieved 0.3821 by using the reference evidence
to select among all five retrieval outcomes. This gap shows that different
questions can benefit from different sizes. It does not establish that
a router can predict those choices for a new question.

The early adaptive methods used the question embedding to predict a size,
or pooled chunks from all sizes into a single ranking. The selected
logistic-regression router reached 0.2853 retrieval F1. Raw mixed-size
retrieval reached 0.2562, and overlap-aware deduplication raised this to
0.2618. Both mixtures remained below fixed 40 and favoured short chunks,
even after deduplication.

\newpage
The next experiments asked whether a small Qwen model could make a better
decision from the question wording. They covered zero-shot prompting,
numeric-output fine-tuning, restricted alias classification, class-balanced
loss, a dedicated classification head, explicit token-count instructions,
and a learning-rate and training-duration search. The strongest
Qwen-only retrieval result was 0.2865 when fine-tuning the model to choose
among five single-token labels with uniform loss weights. Numeric
fine-tuning increased label accuracy while favouring
large chunks and reducing retrieval F1. The contrast exposed a weakness
of the target: total evidence length does not necessarily identify the
best size for each of five retrieved chunks.

The study then added information from the document. Similarity trees
describe how question-to-chunk scores change between aligned coarse and
fine spans. Linear classifiers and boosted trees were evaluated on
level-wise and hierarchical features, followed by fusion with Qwen
features. The initial fusion used Qwen predictions on examples that had
trained the base model. A corrected run used paper-grouped out-of-fold
predictions: each training paper was predicted by a model fitted on other
papers. This removed the known contamination in the training features.
Corrected fusion reached 0.2783 retrieval F1, below the original result.

The strongest learned question-level method used expected-regret
regression. It learned the loss in retrieval F1 associated with each
size, making supervision reflect the consequences of a routing decision.
It reached 0.3075 mean retrieval F1, ahead of the preceding learned
routers. Its difference from fixed
40-token retrieval was $-0.00835$, with a paired paper-cluster bootstrap
95\% interval of $[-0.01266,-0.00416]$. Within the evaluated development
sample, the fixed baseline remained stronger.

Finally, local routing allowed different regions of the same paper to use
different sizes. The router learned from 4,081 question--tree examples
drawn from 2,101 questions with mappable evidence. At inference, it chose one
size within each of the five highest-ranked 160-token trees and returned
one chunk per tree.
This reached 0.3053 F1, against 0.3075 for a fixed-40 method using the same
trees. A frozen Qwen feature describing the anticipated context need as
short, medium, or long was also tested. Free generation sometimes
produced answers instead of a category. Restricting the decision to
the three labels resolved that problem, but almost every question
received \enquote{medium}. Retrieval
F1 was 0.3045, with no demonstrated gain over the local router alone.

The experiments identify retrieval-utility supervision as a promising
direction, without demonstrating an advantage over the corresponding
fixed baselines. Comparisons also depend on the returned text budget:
five chunks can contain very different amounts of text. Since validation
informed repeated development decisions, the next step is evaluation on
untouched data, followed by additional datasets and measurements of
generated-answer quality.
